\chapter{Determinants}

\section{Introduction}

In mathematics, the \textbf{determinant} is a scalar value that is a
function of the entries of a square matrix. The determinant of a
matrix is commonly denoted $\det(A)$, $\det A$, or
$|A|$. Its value characterizes some properties of the matrix and the
linear map represented by the matrix.  In particular, the determinant
is nonzero if and only if the matrix is invertible and the linear map
represented by the matrix is an isomorphism. The determinant of a
product of matrices is the product of their determinants.

The determinant of a $2 \times 2$ matrix is
$$\begin{vmatrix} a & b\\c & d \end{vmatrix}=ad-bc,$$
and the determinant of a $3 \times 3$ matrix is
$$\begin{vmatrix} a & b & c \\ d & e & f \\ g & h & i \end{vmatrix}=
aei + bfg + cdh - ceg - bdi - afh.$$

The determinant of an matrix can be defined in several equivalent
ways, the most common being Leibniz formula, which expresses the
determinant as a sum of $n!$ (the factorial of $n$) signed products of
matrix entries. It can be computed by the Laplace expansion, which
expresses the determinant as a linear combination of determinants of
submatrices, or with Gaussian elimination, which expresses the
determinant as the product of the diagonal entries of a diagonal
matrix that is obtained by a succession of elementary row operations.

Determinants can also be defined by some of their properties. Namely,
the determinant is the unique function defined on the matrices that
has the four following properties:

\begin{enumerate}
\item The determinant of the identity matrix is $1$.
\item The exchange of two rows multiplies the determinant by
  $-1$. \label{row_exchange}
\item Multiplying a row by a number multiplies the determinant by this
  number. \label{row_multiple}
\item Adding to a row a multiple of another row does not change the
  determinant. \label{row_add}
\end{enumerate}

The above properties relating to rows (properties
\ref{row_exchange}-\ref{row_add}) may be replaced by the
corresponding statements with respect to columns.

Determinants occur throughout mathematics. For example, a matrix is
often used to represent the coefficients in a system of linear
equations, and determinants can be used to solve these equations
(Cramer's rule), although other methods of solution are
computationally much more efficient.  Determinants are used for
defining the characteristic polynomial of a matrix, whose roots are
the eigenvalues. In geometry, the signed $n$-dimensional volume of a
$n$-dimensional parallelepiped is expressed by a determinant, and the
determinant of (the matrix of) a linear transformation determines how
the orientation and the $n$-dimensional volume are transformed. This is
used in calculus with exterior differential forms and the Jacobian
determinant, in particular for changes of variables in multiple
integrals.

\section{Two by two matrices}

The determinant of a matrix
\begin{equation}
  \begin{pmatrix} a & b \\c & d \end{pmatrix}
  \label{eq:2by2}
\end{equation}
  is denoted either by ``det''
or by vertical bars around the matrix, and is defined as
$$\det \begin{pmatrix} a & b \\c & d \end{pmatrix} = \begin{vmatrix} a & b \\c & d \end{vmatrix} = ad - bc.$$
For example,
$$\det \begin{pmatrix} 3 & 7 \\1 & -4 \end{pmatrix} = \begin{vmatrix} 3 & 7 \\ 1 & {-4} \end{vmatrix} = (3 \cdot (-4)) - (7 \cdot 1) = -19.$$

\subsection{First properties}

The determinant has several key properties that can be proved by direct
evaluation of the definition for $2 \times 2$-matrices, and that
continue to hold for determinants of larger matrices. They are as
follows: first, the determinant of the identity matrix
$\begin{pmatrix}1 & 0 \\ 0 & 1 \end{pmatrix}$ is 1. Second, the
determinant is zero if two rows are the same:
$$\begin{vmatrix} a & b \\ a & b \end{vmatrix} = ab - ba = 0.$$
This holds similarly if the two columns are the same. Moreover,
$$\begin{vmatrix} a & b + b' \\ c & d + d' \end{vmatrix} = a(d+d')-(b+b')c = \begin{vmatrix}a & b\\ c & d \end{vmatrix} + \begin{vmatrix}a & b' \\ c & d' \end{vmatrix}.$$
Finally, if any column is multiplied by some number $r$ (i.e., all
entries in that column are multiplied by that number), the determinant
is also multiplied by that number:
$$\begin{vmatrix} r \cdot a & b \\ r \cdot c & d \end{vmatrix} = rad - brc = r(ad-bc) = r \cdot \begin{vmatrix} a & b \\c & d \end{vmatrix}.$$

\section{Geometric meaning}

\begin{figure}
\centering
\includegraphics[width=0.5\textwidth]{Area_parallellogram_as_determinant.png}
\caption{\label{fig:geo} The area of the parallelogram is the absolute
  value of the determinant of the matrix formed by the vectors
  representing the parallelogram's sides.}
\end{figure}

If the matrix entries are real numbers, the matrix $A$ can be used to
represent two linear maps: one that maps the standard basis vectors to
the rows of $A$, and one that maps them to the columns of $A$. In
either case, the images of the basis vectors form a parallelogram that
represents the image of the unit square under the mapping. The
parallelogram defined by the rows of the matrix \eqref{eq:2by2} is the
one with vertices at $(0, 0)$, $(a, b)$, $(a+c, b+d)$, and $(c, d)$,
as shown in Figure \ref{fig:geo}.

The absolute value of $ad - bc$ is the area of the parallelogram, and
thus represents the scale factor by which areas are transformed by
$A$. (The parallelogram formed by the columns of $A$ is in general a
different parallelogram, but since the determinant is symmetric with
respect to rows and columns, the area will be the same.)

The absolute value of the determinant together with the sign becomes
the \emph{oriented area} of the parallelogram. The oriented area is
the same as the usual area, except that it is negative when the angle
from the first to the second vector defining the parallelogram turns
in a clockwise direction (which is opposite to the direction one would
get for the identity matrix.

[Original Wikipedia article continues\ldots]
